%!Mode:: "TeX:UTF-8"
%!TEX program = xelatex

% ============================================
% 山东大学学位论文 —— 主文件
%
% 本文件只负责"结构"：文档类选项、前置/正文/后置的顺序。
% 论文信息一律写在 sdusetup.tex，不要写在这里。
% ============================================

% ============================================
% 文档类选项（全部可选；不写则用默认值）
%
%   master / doctor       学位层次        默认 master
%   colorcover / bwcover  封面颜色        默认按学位自动：硕士=山大红、博士=黑白
%   blindreview           盲审模式        默认关（开启后封面/扉页姓名、学号、导师显示「（隐去）」）
%   chichapter            章标题「第一章」 默认开
%   arabicchapter         章标题「1」      默认关（指南三（七）对工学的建议形式）
%   noflyleaf             不出扉页        默认关（按指南二（二）出封面+扉页两页）
%   flushbottom           段距可被拉伸    默认关（默认 raggedbottom：段距恒定、页底参差）
%   fontset=...           字体集          默认 windows（另有 mac / ubuntu / fandol）
%   biblatex              兼容旧文件用    恒开（参考文献固定 biblatex + GB/T 7714-2015）
% ============================================
\documentclass[
  master,
  % bwcover,              % 显式黑白封面（覆盖自动判定）
  % colorcover,           % 显式红色封面（覆盖自动判定）
  % blindreview,          % 盲审：封面与扉页隐去可识别信息
  % arabicchapter,        % 章标题改为「1 绪论」（默认「第一章 绪论」）
  % noflyleaf,            % 不出扉页（只出封面）
  % flushbottom,          % 页底强制对齐（默认段距恒定、页底参差）
  fontset=windows,
]{sduthesis}

% ============================================
% 论文基本信息（唯一必改文件）
% ============================================
%!Mode:: "TeX:UTF-8"
% ============================================================
%  sdusetup.tex —— 论文基本信息
%
%  ★ 这是你唯一"必须"修改的文件。
%  升级模板时，只需保留本文件与 chapters/、figures/、refs.bib，
%  其余文件可被新版整包覆盖，你的论文信息不会丢失。
%
%  ⚠ 花括号内不要留空行：TeX 中空行等于 \par，会打断键值列表。
%     需要分隔时用 % 注释行。
%
%  封面字段与 2024 版《学位论文封面（扉页）》一一对应。
% ============================================================

\sdusetup{
  % ---- 题目 ----
  title       = {论文题目待定},
  titleen     = {Title to Be Determined},
  % ---- 作者与学籍 ----
  author      = {张三},
  studentid   = {2020xxxxxxxx},
  college     = {示例学院},
  trainingmode= {全日制},
  major       = {示例专业},
  supervisor  = {导师姓名},
  cosupervisor= {},
  % ---- 时间与分类 ----
  date        = {2026年6月1日},
  fenlei      = {TK},            % 分类号，按《中国图书资料分类法》填写
  miji        = {公开},          % 公开 / 内部 / 秘密 / 机密
  % ---- 学位类型 ----
  % academic 学术学位 | professional 专业学位 | interdisciplinary 学科交叉
  % equivalent 同等学力 | engineering 工程硕博士培养改革专项
  degreetype  = {academic},
  papertype   = {},              % 仅 engineering 需要，如 专题研究类论文
}




% ============================================
% 参考文献数据库（biblatex + GB/T 7714-2015 恒开）
% ============================================
\addbibresource{thesis.bib}

% 参考文献著录细节（GB/T 7714-2015）：
%   默认 gbmedium=true —— 带 DOI/URL 的条目会输出 [J/OL] 等在线载体标识。
%   若学院要求一律用印刷版标识（[J]/[M]/[C]），取消下一行注释即可：
% \ExecuteBibliographyOptions{gbmedium=false}

% 常用扩展包（模板不预加载，需要时自行加载）：
%   \usepackage{tabularx}            宽表自动换行
%   \usepackage{threeparttable}      表格附注（与 \sdutablenotes 可二选一）
%   \usepackage{siunitx}             单位与数值
%   \usepackage{algorithm} \usepackage{algpseudocode}
%   \usepackage[ruled]{algorithm2e}  ← 算法编号与题注已由模板处理为「算法 3-1」
%   \usepackage{listings}            代码
%   \usepackage{todonotes}           批注
% 提示：hyperref 由模板在导言区加载；cleveref 等依赖它的宏包放在导言区即可。

% ============================================
% 正文开始
% ============================================
\begin{document}

% ===== 封面与声明（无页码、无页眉）=====
\maketitle
\makestatement

% ===== 前置部分（大写罗马数字页码）=====
% 结构顺序依据 2022 指南 二：摘要 → 外文摘要 → 目录 → 符号说明
\clearpage
\pagenumbering{Roman}
\setcounter{page}{1}

% 中文摘要
\begin{cnabstract}
本模板严格依据《山东大学研究生学位论文撰写指南》（2022）与《2024 版学位论文封面（扉页）》编写，
覆盖封面、声明、摘要、目录、正文、参考文献与各类后置页面。

\cnkeywords{学位论文；\LaTeX{}；模板；山东大学}
\end{cnabstract}

% 外文摘要
\begin{enabstract}
This template follows the Shandong University Guidelines for Writing Postgraduate
Degree Theses (2022). It covers the cover page, declaration, abstracts, table of
contents, main matter, bibliography and back matter.

\enkeywords{Thesis; LaTeX; Template; Shandong University}
\end{enabstract}

% 目录（中英文各一份）
\tableofcontents
\tableofengcontents

% 符号说明（可选；按 2022 指南二（七），位置在目录之后、正文之前）
%!Mode:: "TeX:UTF-8"

\makesymbolnotation

本文所用主要符号说明如下：

\begin{denotation}[2.8cm]
  \item[$\rho$]        密度
  \item[$\mathbf{v}$]  速度矢量
  \item[$t$]           时间
  \item[$\nabla$]      梯度算子
  \item[$\mathcal{L}$] 损失函数
  \item[$\theta$]      模型参数
  \item[$\alpha$]      学习率
  \item[$\eta$]        学习率衰减系数
  \item[$\lambda$]     正则化系数
  \item[$\mathbf{x}_t$] 第 $t$ 个时间步的观测向量
  \item[$\mathbf{m}_t$] 第 $t$ 个时间步的缺失掩码
  \item[$\hat{\mathbf{x}}$] 插补后的重构序列
  \item[$\mathbb{R}$]  实数集
  \item[$\mathbb{C}$]  复数集
  \item[$\mathrm{e}$]  自然对数的底
  \item[$\mathrm{i}$]  虚数单位
  \item[$\Delta t$]    采样时间间隔
  \item[$\sigma$]      观测噪声标准差
  \item[$\mu$]         均值
  \item[$\Sigma$]      协方差矩阵
  \item[$\mathrm{SOC}$] 荷电状态（State of Charge）
  \item[$\mathrm{SOH}$] 健康状态（State of Health）
  \item[$\mathrm{RUL}$] 剩余使用寿命（Remaining Useful Life）
  \item[$\mathrm{BMS}$] 电池管理系统（Battery Management System）
  \item[$\mathrm{MAE}$] 平均绝对误差
  \item[$\mathrm{RMSE}$] 均方根误差
  \item[$\mathrm{MRE}$] 平均相对误差
  \item[$\mathrm{MCAR}$] 完全随机缺失
  \item[$\mathrm{MAR}$]  随机缺失
  \item[$\mathrm{MNAR}$] 非随机缺失
  \item[$\mathrm{NCDE}$] 神经受控微分方程
  \item[$\mathrm{ODE}$]  常微分方程
  \item[$\mathrm{RNN}$]  循环神经网络
  \item[$\mathrm{LSTM}$] 长短期记忆网络
  \item[$\mathrm{GAN}$]  生成对抗网络
  \item[$\mathrm{EKF}$]  扩展卡尔曼滤波
  \item[$\mathrm{ECM}$]  等效电路模型
  \item[$\mathrm{SVM}$]  支持向量机
  \item[$\mathrm{PDE}$]  偏微分方程
  \item[$\mathrm{GPU}$]  图形处理器
\end{denotation}

注：符号按字母顺序排列；如有未列出的符号，请按同样格式补充。


% ===== 正文部分（阿拉伯数字页码）=====
\clearpage
\pagenumbering{arabic}
\setcounter{page}{1}

%!Mode:: "TeX:UTF-8"

\chapter{绪论}
\echapter{Introduction}  % 英文目录标题

\section{研究背景与意义}
\esection{Research Background and Significance}

学位论文是研究生培养的重要环节，规范的论文格式对于保证学位授予质量具有重要意义。相关背景与排版规范可参见文献\cite{example2024,knuth1984texbook}。

\section{研究内容}

本文的主要研究内容包括：
\begin{enumerate}
  \item 研究内容一；
  \item 研究内容二；
  \item 研究内容三。
\end{enumerate}

\section{公式与图表约定}

本模板的公式按章编号，编号右端对齐，并与公式之间以点线连接（指南三（十））。例如质量守恒的
微分形式可写为

\begin{sduequation}
  \frac{\partial \rho}{\partial t} + \nabla\cdot(\rho\,\mathbf{v}) = 0
  \label{eq:continuity}
\end{sduequation}

式~\eqref{eq:continuity} 中 $\rho$ 为密度，$\mathbf{v}$ 为速度矢量。多行公式仍可使用
\texttt{align} 等环境，编号同样按章编排。

图与表同样按章编号。需要补充说明时使用图注（排在图题下方），用阿拉伯数字按顺序编排
（指南三（八））：

\begin{figure}[htbp]
  \centering
  \includegraphics[width=0.5\textwidth]{figures/SDUWordsColor.pdf}
  \caption{校名标识（示例图）}
  \label{fig:sdu}
  \begin{sdufigurenotes}
    \item 图注位于图题下方，用阿拉伯数字按顺序编号。
    \item 正式论文中请将本图替换为自绘插图。
  \end{sdufigurenotes}
\end{figure}

\section{论文结构}

本文共分为五章，各章内容安排如下：

第一章为绪论，介绍研究背景和意义。

第二章为相关研究，综述国内外研究现状。

第三章为研究方法，详细介绍本文提出的方法。

第四章为实验，验证方法的有效性。

第五章为总结与展望。

%!Mode:: "TeX:UTF-8"

\chapter{相关研究}
\echapter{Related Work}

\section{国内外研究现状}
\esection{Research Status at Home and Abroad}

本节综述相关领域的研究进展。

\section{现有方法分析}

对现有方法进行分析比较。

%!Mode:: "TeX:UTF-8"

\chapter{研究方法}
\echapter{Methodology}

\section{方法概述}
\esection{Method Overview}

详细介绍本文提出的研究方法。

\section{关键技术}

\subsection{技术一}

技术一的详细描述。

\subsection{技术二}

技术二的详细描述。

%!Mode:: "TeX:UTF-8"

\chapter{实验}
\echapter{Experiments}

\section{实验设置}
\esection{Experimental Setup}

介绍实验环境和数据集。

\section{实验结果}

展示实验结果和分析。

\begin{table}[htbp]
  \centering
  \caption{实验结果对比}
  \begin{tabular}{lcc}
    \toprule
    方法 & 准确率 & 召回率 \\
    \midrule
    方法A & 0.85 & 0.82 \\
    方法B & 0.88 & 0.85 \\
    本文方法 & \textbf{0.92} & \textbf{0.90} \\
    \bottomrule
  \end{tabular}
  % 表附注：排在表的下方，用阿拉伯数字按顺序编排（指南三（九））
  \begin{sdutablenotes}
    \item 表中数值为 10 次重复实验的均值。
    \item 加粗表示该列最优结果。
  \end{sdutablenotes}
\end{table}

\section{长表格与续表}

当表格需要转页接排时，续页应重复表号（写作「续表 4-2」）并重复表头（指南三（九））。
本模板用 \texttt{longtable} 配合 \verb|\sdutablecaption| 与 \verb|\sdutablecont| 实现，
下表即是一个跨页示例。

\begin{longtable}{ccc}
  \sdutablecaption{跨页长表示例}\label{tab:long}\\
  \toprule
  序号 & 测量值 & 偏差 \\
  \midrule
  \endfirsthead
  % 续页：重复「续表 x-y + 表题」，并重复表头
  \sdutablecont\\
  \toprule
  序号 & 测量值 & 偏差 \\
  \midrule
  \endhead
  \bottomrule
  \endfoot
  1 & 0.312 & 0.004 \\
  2 & 0.318 & 0.003 \\
  3 & 0.309 & 0.006 \\
  4 & 0.321 & 0.002 \\
  5 & 0.315 & 0.005 \\
  6 & 0.311 & 0.004 \\
  7 & 0.317 & 0.003 \\
  8 & 0.313 & 0.005 \\
  9 & 0.319 & 0.002 \\
  10 & 0.310 & 0.006 \\
  11 & 0.316 & 0.004 \\
  12 & 0.314 & 0.005 \\
  13 & 0.320 & 0.003 \\
  14 & 0.308 & 0.007 \\
  15 & 0.322 & 0.002 \\
  16 & 0.315 & 0.004 \\
  17 & 0.311 & 0.005 \\
  18 & 0.317 & 0.003 \\
  19 & 0.313 & 0.006 \\
  20 & 0.319 & 0.002 \\
  21 & 0.310 & 0.005 \\
  22 & 0.316 & 0.004 \\
  23 & 0.314 & 0.006 \\
  24 & 0.320 & 0.003 \\
  25 & 0.308 & 0.007 \\
  26 & 0.322 & 0.002 \\
  27 & 0.315 & 0.004 \\
  28 & 0.311 & 0.005 \\
  29 & 0.317 & 0.003 \\
  30 & 0.313 & 0.006 \\
  31 & 0.319 & 0.002 \\
  32 & 0.310 & 0.005 \\
  33 & 0.316 & 0.004 \\
  34 & 0.314 & 0.006 \\
  35 & 0.320 & 0.003 \\
  36 & 0.308 & 0.007 \\
  37 & 0.322 & 0.002 \\
  38 & 0.315 & 0.004 \\
  39 & 0.311 & 0.005 \\
  40 & 0.317 & 0.003 \\
\end{longtable}

%!Mode:: "TeX:UTF-8"

\chapter{总结与展望}
\echapter{Conclusions and Future Work}

\section{工作总结}
\esection{Work Summary}

本文主要工作总结如下：
\begin{enumerate}
  \item 完成了工作一；
  \item 完成了工作二；
  \item 完成了工作三。
\end{enumerate}

\section{未来展望}

未来研究工作包括：
\begin{enumerate}
  \item 研究方向一；
  \item 研究方向二。
\end{enumerate}


% ===== 后置部分 =====

% 参考文献（biblatex + GB/T 7714-2015；标题进目录、另起一页）
\printbibliography[heading=bibintoc,title={参考文献}]

% 致谢
%!Mode:: "TeX:UTF-8"

% 致谢部分自动添加到目录，无需手动添加英文目录
\acknowledgement

三年的研究生生活即将结束，感谢所有给予我帮助和支持的人们。

首先感谢我的导师，在学术研究和论文撰写过程中给予了悉心指导。

感谢实验室的同学和朋友们的帮助和支持。

最后感谢我的家人，是你们的理解和支持让我能够顺利完成学业。


% 攻读学位期间发表的学术论文目录（可选）
% %!Mode:: "TeX:UTF-8"

\makepublications

\begin{enumerate}[label={[}\arabic*{]}]
\item 作者姓名. 论文标题[J]. 期刊名称, 2024, 卷(期): 起止页码.（SCI/EI收录）

\item 作者姓名. 论文标题[C]. 会议名称, 会议地点, 2024: 起止页码.（EI收录）

\item 作者姓名. 专利名称[P]. 专利号, 授权日期.
\end{enumerate}

注：请按发表时间顺序列出所有论文，注明收录情况和本人排序。


% 学位论文评阅及答辩情况（可选）
% %!Mode:: "TeX:UTF-8"

\makedefense

\begin{center}
\zihao{4}\bfseries
学位论文评阅及答辩情况
\end{center}

\vspace{1cm}

\zihao{-4}

{\bfseries 一、论文评阅人}

评阅人姓名：\hrulefill

职称/职务：\hrulefill

工作单位：\hrulefill

评阅意见：\hrulefill

\vspace{0.5cm}

评阅人姓名：\hrulefill

职称/职务：\hrulefill

工作单位：\hrulefill

评阅意见：\hrulefill

\vspace{1cm}

{\bfseries 二、答辩委员会成员}

\begin{tabular}{|c|c|c|c|}
\hline
姓名 & 职称 & 工作单位 & 备注 \\
\hline
& & & 主席 \\
\hline
& & & 委员 \\
\hline
& & & 委员 \\
\hline
& & & 委员 \\
\hline
& & & 委员 \\
\hline
\end{tabular}

\vspace{1cm}

{\bfseries 三、答辩情况}

答辩日期：\hrulefill

答辩地点：\hrulefill

答辩决议：\hrulefill

\vspace{2cm}

记录人签名：\hrulefill\quad 日期：\hrulefill


% 附录（可选）
% \appendix
% \input{chapters/appendix}

\end{document}



